\SourceLeaf{491}%
Skagen angestellt wurden, könnte es dem Vertrauen schaden, das die Ergebnisse für sich beanspruchen sollen, wenn jene Beobachtungen mit verschiedenen Instrumenten und sogar mit Instrumenten ganz verschiedener Art gewonnen worden wären. Am besten wird es daher sein, wenn die hiesigen Beobachtungen sowohl am Meridiankreis als auch am Zenitsektor angestellt werden, und zwar an \textit{demselben} Zenitsektor, mit dem an allen anderen Orten beobachtet wurde.

Der Zenitsektor, mit dem die dänischen Bestimmungen in Lauenburg, Lysabbel und Skagen sowie in Kopenhagen in den Jahren 1819 und 1820 vorgenommen wurden, ist eben derselbe, mit dem der verstorbene General \textsc{Mudge} bei der englischen Gradmessung im Jahr 1802 beobachtet hat. Mit eben diesem Instrument wurde im Jahr 1818 von englischen und französischen Gelehrten gemeinsam in Dünkirchen, dem nördlichsten Endpunkt der französischen Gradmessung, eine Reihe von Beobachtungen angestellt. Danach lieh es die englische Regierung für die dänische Gradmessung für 2 Jahre aus. All diese Umstände machen es umso wünschenswerter, dass dieses Instrument auch bei der hannoverschen Gradmessung verwendet wird. Dabei wird es besonders wichtig und lehrreich sein, es durch gleichzeitige und an einem einzigen Ort angestellte Beobachtungen mit dem hiesigen Reichenbachschen Meridiankreis zu vergleichen.

Es besteht kein Zweifel, dass dieses Instrument für die hannoversche Gradmessung mit derselben Bereitwilligkeit ausgeliehen wird wie für die dänische. — — — Ich füge noch folgende Umstände hinzu. Dieser von \textsc{Ramsdon} gefertigte Zenitsektor gehört dem Board of Ordnance und untersteht in letzter Instanz dem Herzog von \textsc{Wellington}, mittelbar aber nach dem vor Kurzem eingetretenen Tod des Generals \textsc{Mudge} dem Oberstleutnant \textsc{Colby}.

Aus der vorstehenden Darstellung geht hervor, dass sich der Zeitpunkt, zu dem die größeren, besondere Vorkehrungen erfordernden Arbeiten der Gradmessung beginnen können, im Augenblick noch nicht bestimmt angeben lässt. In jedem Fall wird im Lauf des gegenwärtigen Jahres nur ein Teil der astronomischen Beobachtungen an der hiesigen Sternwarte (mit denen ich bereits angefangen habe) und höchstens einige Vorarbeiten in der hiesigen Umgebung geleistet werden können.

Göttingen, den 11. August 1820.
\begin{flushright}C. F. \textsc{Gauss}.\end{flushright}

\par\bigskip
\pdfbookmark[0]{P. M. über die Anwesenheit des Hofrats Gauss bei einem Teil der dänischen Gradmessung.}{pm-gradmessung}
\begin{center}\large P. M. ÜBER DIE ANWESENHEIT DES HOFRATS GAUSS BEI EINEM TEIL DER ARBEITEN DER DÄNISCHEN GRADMESSUNG.\end{center}
\par
Über die im königlich dänischen Gebiet angeordnete und dem dänischen Astronomen \textsc{Schumacher} übertragene Gradmessung, ihren Zweck, die Vorteile, die man für die Wissenschaften davon erwarten darf, sowie über die dabei eingesetzten Hilfsmittel hatte ich schon früher die Ehre, Ew. — — in aller Ergebenheit zu berichten. Es freut mich umso mehr, dass gerade jetzt ein anderer Astronom, Doktor \textsc{Olbers}, der im vergangenen Jahr ebenfalls Augenzeuge der Arbeiten in Lauenburg war, in der Zeitschrift Allgemeine Geographische Ephemeriden (Bd. 17, Heft 3) Ansichten über dieses Unternehmen veröffentlicht hat, die ganz mit meinen übereinstimmen. Mein eigenes Urteil über die Person des dänischen Astronomen könnte nämlich insofern vielleicht befangen erscheinen, als dieser früher (1809) seine Ausbildung an der hiesigen Universität abgeschlossen hatte.
